\section{Implementation}

% Structure du code et choix techniques.

\subsection{Organisation du code}

Le moteur sépare trois opérations : calcul du pas libre, projection des positions
sur les contraintes, puis relecture de la vitesse par déplacement réalisé.
Les modules de modèle fournissent la direction souhaitée et les propriétés du milieu.
Le scénario assemble ces fonctions avec la géométrie, les rayons et la graine.
Les observables reçoivent ensuite des tableaux de trajectoire, sans dépendre du simulateur.

Dans le banc de couloir, la direction est constante vers l'aval. Une carte par bandes
associe à chaque centre d'agent son milieu ; une frontière appartient à la bande de droite.
Le milieu est évalué au début du pas : un changement de zone influence donc le pas suivant.
La projection d'exclusion impose les contacts entre disques et les trois parois du couloir.
La section de mesure, fournie par la même géométrie, ne constitue pas une paroi.

\subsection{Recherche des voisins}

% Grille de hachage spatial, pour eviter le cout en $N^2$.

\subsection{Reproductibilite}

La chaîne disponible est
\[
\text{configuration YAML} \longrightarrow \text{simulation}
\longrightarrow \text{archive NPZ} \longrightarrow \text{mesures et relecture}.
\]
Les configurations sont versionnées. Les clés inconnues ou dupliquées et les valeurs
physiques invalides sont refusées. Tous les paramètres du couloir sont obligatoires ;
son initialisation utilise exclusivement le générateur commun à graine explicite.
Les contrôles imposent notamment $\dt\leq\min_m\tau_m$ et un déplacement souhaité
puis effectivement mesuré inférieur ou égal à $r/5$ par pas. Ces garde-fous ne remplacent
pas une étude de convergence. Une non-convergence de la projection interdit l'archivage.

L'archive conserve les instants, positions, vitesses et rayons en unités SI, ainsi que
la configuration effective, la graine, les versions logicielles et les diagnostics.
Les métadonnées utilisent JSON et les tableaux se chargent sans pickle. Les tests
comparent exactement les tableaux relus et ceux de deux calculs à graine identique
dans un même environnement ; cette propriété ne garantit pas l'identité entre versions
de dépendances. Le fichier existant n'est jamais écrasé.

La commande \texttt{rejouer} anime le couloir uniquement à partir de l'archive : positions,
rayons, horloge et compteur de passages proviennent des données enregistrées. Le curseur
choisit le dernier instant archivé inférieur ou égal au temps demandé ; le GIF inclut
le dernier instant. Lire/Pause, vitesse de lecture et export PNG/GIF ne modifient pas
la dynamique. Les archives gardent tous les agents et tous les pas en mémoire : elles
ne constituent ni une reprise de calcul ni un état de sortie absorbant.
